\documentclass[12pt, a4paper]{article}
\usepackage{../../../myconfig} % Gọi file cấu hình từ ngoài gốc vào

% ==========================================
% BẮT ĐẦU NỘI DUNG TÀI LIỆU
% ==========================================
\begin{document}

% Truyền 3 tham số: {Nhãn cam} {Chữ to chính giữa} {Chữ ở góc phải Header}
\titlebox{Chủ đề: Hàm số}{Cấp số cộng- cấp số nhân}{Hàm số: CSC-CSN}

\drawdashlines{20}
\newpage


% --- CÂU 1 ---
\questionLinesManual{5}{1}{Trong các dãy số sau, dãy số nào là một cấp số cộng?}{
    \begin{tasks}(4)
        \task $1; -2; -4; -6; -8$
        \task $1; -3; -6; -9; -12$
        \task $1; -3; -7; -11; -15$
        \task $1; -3; -5; -7; -9$
    \end{tasks}
}



% --- CÂU 2 ---
\questionLinesManual{5}{2}{Cho một cấp số cộng \( (u_n) \) có \( u_1 = \dfrac{1}{3} \), \( u_8 = 26 \). Tìm công sai \( d \)}{
    \begin{tasks}(4)
        \task \( d = \dfrac{11}{3} \)
        \task \( d = \dfrac{10}{3} \)
        \task \( d = \dfrac{3}{10} \)
        \task \( d = \dfrac{3}{11} \)
    \end{tasks}
}

% --- CÂU 3 ---
\questionLinesManual{5}{3}{Cho cấp số cộng \( (u_n) \) có \( u_1 = -2 \) và công sai \( d = 3 \). Tìm số hạng \( u_{10} \).}{
    \begin{tasks}(4)
        \task \( u_{10} = -2 \cdot 3^9 \)
        \task \( u_{10} = 25 \)
        \task \( u_{10} = 28 \)
        \task \( u_{10} = -29 \)
    \end{tasks}
}

% --- CÂU 30 ---
\questionLinesManual{5}{4}{Cho cấp số cộng \( (u_n) \), biết \( u_2 = 3 \) và \( u_4 = 7 \). Giá trị của \( u_{3} \) bằng}{
    \begin{tasks}(4)
        \task $27$
        \task $31$
        \task $35$
        \task $29$
    \end{tasks}
}


% --- CÂU 4 ---
\questionLinesManual{5}{5}{Cho cấp số cộng \( (u_n) \) có \( u_1 = 1 \) và công sai \( d = 2 \). Tổng \( S_{10} = u_1 + u_2 + u_3 + \ldots + u_{10} \) bằng:}{
    \begin{tasks}(4)
        \task \( S_{10} = 110 \)
        \task \( S_{10} = 100 \)
        \task \( S_{10} = 21 \)
        \task \( S_{10} = 19 \)
    \end{tasks}
}

% --- CÂU 5 ---
\questionLinesManual{5}{6}{Dãy số nào sau đây \textbf{không} phải là cấp số nhân?}{
    \begin{tasks}(4)
        \task $1; 2; 4; 8; \ldots$
        \task $3; 3^2; 3^3; 3^4; \ldots$
        \task $4; 2; \dfrac{1}{2}; \dfrac{1}{4}; \ldots$
        \task $\dfrac{1}{\pi}; \dfrac{1}{\pi^2}; \dfrac{1}{\pi^4}; \dfrac{1}{\pi^6}; \ldots$
    \end{tasks}
}

% --- CÂU 6 ---
\questionLinesManual{5}{7}{Cho cấp số nhân \( (u_n) \) có \( u_1 = -3 \) và \( q = \dfrac{2}{3} \). Mệnh đề nào sau đây đúng?}{
    \begin{tasks}(4)
        \task \( u_5 = -\dfrac{27}{16} \)
        \task \( u_5 = -\dfrac{16}{27} \)
        \task \( u_5 = \dfrac{16}{27} \)
        \task \( u_5 = \dfrac{27}{16} \)
    \end{tasks}
}

% --- CÂU 7 ---
\questionLinesManual{5}{8}{Một cấp số nhân có 6 số hạng, số hạng đầu bằng 2 và số hạng thứ sáu bằng 486. Tìm công bội \( q \) của cấp số nhân đã cho.}{
    \begin{tasks}(4)
        \task \( q = 3 \)
        \task \( q = -3 \)
        \task \( q = 2 \)
        \task \( q = -2 \)
    \end{tasks}
}

% --- CÂU 31 ---
\questionLinesManual{5}{9}{Với giá trị \( x \) nào dưới đây thì các số \(-4; x; -9\) theo thứ tự đó lập thành một cấp số nhân?}{
    \begin{tasks}(4)
        \task \( x = 36 \)
        \task \( x = -\dfrac{13}{2} \)
        \task \( x = 6 \)
        \task \( x = -36 \)
    \end{tasks}
}

% --- CÂU 8 ---
\questionLinesManual{5}{10}{Cho cấp số nhân \( (u_n) \) có \( u_1 = -3 \) và \( q = -2 \). Tính tổng 10 số hạng đầu tiên của cấp số nhân đã cho.}{
    \begin{tasks}(4)
        \task \( S_{10} = -511 \)
        \task \( S_{10} = -1025 \)
        \task \( S_{10} = 1025 \)
        \task \( S_{10} = 1023 \)
    \end{tasks}
}



\end{document}